In my weld-seam HMI, modified-DH kinematics linked joint controls to an OpenGL digital twin; in the separate CR5 project, workspace checks and execution timing connected software decisions to physical motion. I can integrate these layers, but I need a more systematic pathway from robot mechanics through computer control and mechatronics to project and industry-based validation. NUS's MSc in Mechanical Engineering would help me develop that pathway for intelligent machines.

If the current offerings remain available in my intake, ME5402 Advanced Robotics and ME5422 Computer Control and Applications would strengthen the robotics and control foundations beneath my integration decisions. Subject to availability and approval, the optional four-unit ME5001A Mechanical Engineering Project and optional four-unit ME5888M Mechanical Engineering Internship could form a four-plus-four combination within the eight-unit cap on project and internship courses. ME5001A could let me formulate a bounded Control \& Mechatronics or Manufacturing question and investigate it through a defined validation method. ME5888M at a local industrial company could let me apply and refine that method under industrial constraints, where feasibility and integration matter alongside laboratory performance. I would bring experience making coordinate, workspace, and timing assumptions testable across a digital model and real robot. In robotics or intelligent-manufacturing R\&D, I intend to develop machines whose mechanical, control, and software behavior can be validated together.
